% 19b-outside.tex: the last ring. Added 2026-09-25 after a cold read found the paper's climax,
% the first software from outside, told early (Section 6) and buried under its own caveats.
\section{The last ring: other people's software}
\label{sec:outside}

\motif{A system that has checked itself is not a system that anyone else has checked.}

\analogy{Every wall, tower and ledger in this town was built by the people who live in it, and
checked by them. One day travellers arrived who had built none of it. They carried their own seals,
spoke the common language of the road, and owed the town nothing. What they did at the gate is the
only report the town did not write about itself.}

Every ring in this map exists because the ring inside it could be correct and still unable to prove
what a society needs to know. The last ring is the one no amount of internal work can supply: software
written by someone else, which does not share the author's assumptions about what a correct artifact
contains, meeting Polaris at one of its doors. The door is the OpenID4VP verifier of
Section~\ref{sec:verifiers}, chosen because it speaks the protocol an outside wallet actually speaks.

\figfit{0.56}{outside-witnesses}{The three things outside the repository that have exercised it, and
the two controls without which the first would mean nothing. A wallet nobody here wrote presented a
credential and was accepted; the same wallet, with the verifier trusting a different issuer key, was
refused; and a re-presentation against an answered request was refused. The Foundation's hosted
suite scored seven refusals itself and left four acceptances for a human reviewer; a later run of
the same plan was certified. A published corpus agrees with both signature witnesses on the
primitive.}{fig:outside-witnesses}

\subsection{The first stranger: a conformance service}

The OpenID Foundation runs a hosted conformance suite with a test plan for exactly this role, in
which the suite itself plays the wallet and signs the credential. It ran all eleven modules of that
plan against the verifier across the open internet, fetching the request object and posting the
response from outside. The seven negative modules, in which the suite sends a presentation broken in
one specific way, are scored by the service itself on the verifier's refusal, which makes them the
closest thing to an external adversary this project has been offered. The four positive modules end
in a state the suite reserves for a human to attest a successful verification. A later run of the same
plan was submitted for certification, and the Foundation lists \code{polaris-oid4vp 1.0.0rc7} as
OpenID Certified to the OpenID4VP 1.0 + HAIP 1.0 Verifier profile. Each run carries two negative
controls, one per direction: a verifier patched to refuse everything is noticed by the positive
modules, and one patched to accept everything is noticed by all seven negatives. \sWit

\subsection{The second stranger: a wallet}

Then a wallet arrived. An unmodified walt.id Wallet API, the published image at a pinned digest,
fetched the signed request object, verified it against the registered trust anchor, matched its
credential to the query, signed a key-binding token with a key it generated and this repository has
never held, encrypted the response, and sent it. The verifier checked the whole chain and accepted it.

It refused Polaris first, and it was right.

The verifier's own key generator had produced a request-signing certificate with no key-usage
extension, so the certificate whose entire job is signing request objects was not marked as usable
for signing. Every observer inside the repository had passed over it: the conformance modules, the
package's tests, the mutation drill over every refusal, every check in the invariant layer. Each of
them inherits the author's model of what such a certificate must contain. The first observer that did
not inherit it found the gap at once. The certificate was fixed, with a test that fails on a leaf
without the extension, and the exchange carries the same two controls as the conformance run, because a
verifier that accepts everything prints the same success line. Anyone can repeat the exchange from a
clean machine in about ten minutes without cloning the repository. \sWit

This is the argument of the whole map in one event. The rings are the author's; the last ring is
everyone else's, and it is the only one that can see what the others share.

\subsection{What the last ring does not yet establish}

Stated once, and in full. The certification is a self-certification that the Foundation reviewed and
published, for one package, one version and one role: it is not an endorsement, not an audit, and not
a statement about the rest of Polaris. The wallet exchange covers one wallet, one credential format,
one presentation path and one classical signature algorithm, which is what the profile pins; it says
nothing about any other wallet, about mobile documents, about the post-quantum path, or about
interoperability in general. The native protocol, the \code{polaris-*/1} formats the rest of this map
is built on, has not been implemented by anyone outside from the documents alone. No operator other
than the author has run the system, and no independent party has audited it. Each of those is a row
on the project's scoreboard, and each row is empty until a named party fills it
(Section~\ref{sec:limitations}). \sExt
